% Generated draft: requires mathematical and source review.
\documentclass[11pt]{article}
\usepackage[margin=25mm]{geometry}
\usepackage{fontspec}
\setmainfont{Latin Modern Roman}
\usepackage{amsmath,amssymb,mathtools,unicode-math}
\setmathfont{Latin Modern Math}
\usepackage{enumitem,needspace,xurl,hyperref}
\setlength{\parindent}{0pt}
\setlength{\parskip}{5pt}
\setlength{\emergencystretch}{3em}
\newcommand{\N}{\mathbb{N}}
\newcommand{\Z}{\mathbb{Z}}
\newcommand{\Q}{\mathbb{Q}}
\newcommand{\R}{\mathbb{R}}
\newcommand{\C}{\mathbb{C}}
\newcommand{\F}{\mathbb{F}}
\newcommand{\A}{\mathbb{A}}
\newcommand{\PP}{\mathbb{P}}
\newcommand{\E}{\mathbb{E}}
\newcommand{\eps}{\varepsilon}
\newcommand{\dd}{\,\mathrm{d}}
\begin{document}
\section*{UnsolvedMath ID: 700813}
\textbf{Source problem code:} catalogue entry
\subsection{Definitions and mathematical statement}
All groups below are finite. A formation is a nonempty class of groups closed under isomorphisms, homomorphic images, and finite subdirect products. For a formation \(\mathcal H\), define the \(\mathcal H\)-residual of \(G\) by
\[
G^{\mathcal H}=\bigcap\{N\trianglelefteq G:G/N\in\mathcal H\}.
\]
Use the standard residual convention for the product of formations:
\[
\mathcal A\mathcal B=\{G:G^{\mathcal B}\in\mathcal A\}.
\]
A chief factor of \(G\) is a quotient \(H/K\), where \(K<H\) are normal subgroups of \(G\) and no normal subgroup of \(G\) lies strictly between them. Put
\[
C_G(H/K)=\{g\in G:ghg^{-1}K=hK\text{ for every }h\in H\}.
\]
A formation \(\mathcal F\) is local if there is an assignment \(f\) taking each prime \(p\) to a formation or the empty class such that
\[
G\in\mathcal F
\quad\Longleftrightarrow\quad
G/C_G(H/K)\in f(p)
\quad\text{for every chief factor }H/K\text{ of }G
\text{ and every prime }p\mid |H/K|.
\]
A formation is non-local if it is not local.

For a formation \(\mathcal F\), write \(D_{\mathrm{nt}}(\mathcal F)\) for the ordered pairs \((\mathcal A,\mathcal B)\) of formations satisfying \(\mathcal F=\mathcal A\mathcal B\) for which this decomposition is non-trivial in the intended source convention. The supplied record does not define that convention; consequently this part of the formulation remains unresolved.

The question attributed to L. A. Shemetkov is whether there exists a local formation \(\mathcal F\) such that
\[
D_{\mathrm{nt}}(\mathcal F)\ne\varnothing
\quad\text{and}\quad
\forall(\mathcal A,\mathcal B)\in D_{\mathrm{nt}}(\mathcal F),\qquad
\mathcal A\text{ and }\mathcal B\text{ are both non-local}.
\]
Here ``in any such a decomposition'' is read as referring to non-trivial decompositions. A requirement covering every decomposition without qualification would include \(\mathcal F=\{1\}\mathcal F=\mathcal F\{1\}\), whose factors are local. The exact intended exclusion of trivial decompositions requires source clarification.

The source additionally states that local products of non-local formations exist. This is retained as a supplied background assertion, without an independent verification or proof.
\subsection{Short English statement}
Does there exist a local formation of finite groups that admits a non-trivial product decomposition, but for which every non-trivial product decomposition has two non-local factors? The source does not specify exactly which decompositions count as non-trivial. It also states that some local formations are products of two non-local formations.
\subsection{Sources}
\begin{itemize}
\item Local UnsolvedMath registry, supplied complete source record: permanent ID \(700813\), exact source code \texttt{catalogue entry}, title \emph{Kourovka Notebook Problem 11.121}. The registry labels the problem open; that label is supplied metadata, not an independently verified mathematical assessment.
\item \emph{The Kourovka Notebook (Unsolved Problems in Group Theory)}, original problem 11.121, attributed to L. A. Shemetkov, 11th issue (1990). The supplied background also identifies the 21st issue edition (2026), \texttt{arXiv:1401.0300v48}. These bibliographic details are supplied attribution; the cited text has not been independently consulted.
\item The supplied collection description names E. I. Khukhro and V. D. Mazurov as editors of the Kourovka Notebook collection covering Issues 1--21 (1965--2026).
\end{itemize}
\end{document}
